\chapter{Expanding smoothings of planar corners}\label{sec:catalogue-expanding-corners}
The convention throughout these entries is $\vec H=X^\perp/2$ for the time-one profile. Thus $M_t=\sqrt t\,M$, $t>0$, is the corresponding flow. The positive weight for variational statements is $e^{|X|^2/4}$; minimizing this weighted area locally will be called $E$-minimizing.

\begin{example}[Even graphical corner expanders]\label{ex:catalogue-expanding-corners}
For every $\kappa>0$ there is a smooth even function $\phi_\kappa:\mathbb R\to\mathbb R$ satisfying
\[
 \frac{\phi_\kappa''}{1+(\phi_\kappa')^2}
 +\frac{x\phi_\kappa'-\phi_\kappa}{2}=0,
 \qquad \phi_\kappa'(0)=0,
 \qquad \lim_{|x|\to\infty}\frac{\phi_\kappa(x)}{|x|}=\kappa.
\]
The profile $\Gamma_\kappa=\{(x,\phi_\kappa(x)):x\in\mathbb R\}$ is a smooth complete properly embedded noncompact curve in $\mathbb R^2$, with topology $\mathbb R$, reflection symmetry, and two ends asymptotic to the rays of $y=\kappa|x|$. It is $E$-minimizing against compact replacements. Its curvature vector is nonzero and points along the upward graph normal: writing $q=\phi_\kappa-x\phi_\kappa'$ gives $q'=-x(1+(\phi_\kappa')^2)q/2$ and $\phi_\kappa''=(1+(\phi_\kappa')^2)q/2$, with $q(0)>0$.

The source gives the quantitative estimate, for $x>0$,
\[
 0\leq\phi_\kappa(x)-\kappa x
 \leq\min\left\{\frac{2e(\kappa+2)}{x}e^{-x^2/4},\,3\kappa+2\right\}.
\]
The flow $\sqrt t\,\Gamma_\kappa$, $t\in(0,\infty)$, smooths the corner at time zero. These are ODE-defined profiles, not elementary closed formulas. Source: \cite{CatDing20}, Section 4, equations (4.9)--(4.15), printed pp.~373--374, supplied PDF pp.~15--16. The figure shows a numerical ODE profile with $\phi(0)=1$. Taking the product with $\mathbb R$ gives an example in $\mathbb R^3$ with one zero principal curvature.
\end{example}

For the planar classification, including the stationary-line endpoint, see
\cite{CatHalldorsson2012}, Theorem~6.1, PDF pp.~13--15. Rotation together
with expansion gives the different family in
Example~\ref{ex:catalogue-rotating-expanding-curves}.

\chapter{Rotational graphical expanders over round cones}\label{sec:catalogue-rotational-expanders}
\begin{example}[Entire radial cone smoothings]\label{ex:catalogue-rotational-expanders}
For $n\geq2$ and $\kappa>0$, consider
\[
 M_\kappa=\{(x,\phi_\kappa(|x|)):x\in\mathbb R^n\}\subset\mathbb R^{n+1},
\]
where $\phi_\kappa$ is the smooth radial solution constructed with
\[
 \frac{\phi_\kappa''}{1+(\phi_\kappa')^2}
 +\frac{n-1}{r}\phi_\kappa'
 +\frac{r\phi_\kappa'-\phi_\kappa}{2}=0,
 \qquad \phi_\kappa'(0)=0,
 \qquad \lim_{r\to\infty}\frac{\phi_\kappa(r)}r=\kappa.
\]
The equation at $r=0$ is interpreted by smooth radial extension; it implies $\phi_\kappa''(0)=\phi_\kappa(0)/(2n)$. This is an ODE construction, not a closed formula.

The hypersurface is smooth, complete, properly embedded, noncompact, rotationally symmetric, diffeomorphic to $\mathbb R^n$, and $E$-minimizing against compact replacements. It has one conical end and nondecreasing height. The estimate
\[
 0\leq\phi_\kappa(r)-\kappa r
 \leq((n+1)\kappa+2)\min\{1,r^{-1}\},\qquad r>0,
\]
records the rate at which it approaches its cone. Its flow is $M_t=\sqrt t\,M_\kappa$ for $t>0$, with initial cone $z=\kappa|x|$. Source: \cite{CatDing20}, Section 4, equations (4.1), (4.7), and (4.8), printed pp.~371--373, supplied PDF pp.~13--15. The interactive figure shows a numerical rotational solution in $\mathbb R^3$ with $\phi(0)=1$.
\end{example}

\chapter{Two-sheeted expanders of rotational double cones}\label{sec:catalogue-two-sheet-expanders}
\begin{example}[Disconnected double-cone smoothings]\label{ex:catalogue-two-sheet-expanders}
Let $n\geq2$, write $(x,z)\in\mathbb R^n\times\mathbb R$, and let
\[
 D_\alpha=\{(x,z):|x|=|z|\tan\alpha\},\qquad 0<\alpha<\pi/2.
\]
With $\phi_{\cot\alpha}$ from Example \ref{ex:catalogue-rotational-expanders}, the two-sheeted expander is
\[
 M_\alpha^{(2)}
 =\{(x,\phi_{\cot\alpha}(|x|)):x\in\mathbb R^n\}
 \mathbin{\cup}
 \{(x,-\phi_{\cot\alpha}(|x|)):x\in\mathbb R^n\}.
\]
This display specifies the surface using ODE-defined functions. Its two components are smooth, complete, properly embedded convex entire graphs, each diffeomorphic to $\mathbb R^n$. The union is noncompact, invariant under rotations about the $z$-axis and under $z\mapsto-z$, and has two oppositely directed conical ends. The flow $\sqrt t\,M_\alpha^{(2)}$, $t>0$, approaches $D_\alpha$ as $t\downarrow0$.

Existence and uniqueness of each rotational single-lobe graphical profile are established in \cite{CatAIC95}, Lemma 2(i), printed p.~1947, supplied PDF p.~12. Theorem 4, printed pp.~1950--1951, PDF pp.~15--16, provides a critical angle $\alpha_{\rm crit}(n+1)\in(0,\pi/2)$. For $\alpha<\alpha_{\rm crit}(n+1)$ the level-set flow of the double cone is exactly the two-sheeted flow. For angles above this threshold there are also connected smoothings, so the graphical uniqueness statement does not assert uniqueness among all flows from the cone. The critical-angle decimal estimates later in that paper are numerical. The $n=2$ figure rotates two reflected numerical profiles; the cone is their asymptotic limit.
\end{example}

\chapter{One-sheeted expanders of rotational double cones}\label{sec:catalogue-one-sheet-expanders}
\begin{example}[Connected double-cone smoothings]\label{ex:catalogue-one-sheet-expanders}
For $n\geq2$ and a neck parameter $a>0$, the initial value problem
\[
 u''=(1+(u')^2)\left(\frac{n-1}{u}+\frac{u-zu'}2\right),
 \qquad u(0)=a,\quad u'(0)=0,
\]
defines a positive, smooth, even profile on $\mathbb R$. Rotating it gives
\[
 M_a=\{(u(z)\omega,z):z\in\mathbb R,\ \omega\in S^{n-1}\}
 \subset\mathbb R^{n+1}.
\]
This is a smooth complete properly embedded noncompact expander with topology $S^{n-1}\times\mathbb R$, rotational symmetry, reflection symmetry, and two conical ends. The profile is increasing for $z>0$. Its limiting aperture depends on $a$; the formula is an ODE specification, not an elementary expression. The evolution is $\sqrt t\,M_a$ for $t>0$.

For every $\alpha>\alpha_{\rm crit}(n+1)$ there exists such a connected expander asymptotic to $D_\alpha$ from Example \ref{ex:catalogue-two-sheet-expanders}. It gives a second thin evolution from that cone alongside the disconnected two-sheeted solution. The cited theorem uses a strict inequality and does not establish a universal uniqueness assertion for the one-sheeted profile. Source: \cite{CatAIC95}, Lemma 2(ii), Lemma 3, and Theorem 4, printed pp.~1947--1951, supplied PDF pp.~12--16. These qualitative expander statements are proved in the source; its separate computed higher-genus shrinker and numerical threshold estimates are not needed here. The interactive figure shows the numerical solution in $\mathbb R^3$ with $a=0.8$ on a finite axial interval.
\end{example}

\chapter{One-sided expanders from mean-convex cones}\label{sec:catalogue-mean-convex-expanders}
\begin{example}[Ding's smooth minimizing expanders]\label{ex:catalogue-mean-convex-expanders}
Let $C\subset\mathbb R^{n+1}$ be an $n$-dimensional $C^{3,\beta}$-regular cone, $0<\beta<1$, regular away from its vertex, with a $C^{3,\beta}$ embedded link in $S^n$. Assume that $C$ is mean convex but not area minimizing, that $\partial\Omega=C$, and that its mean curvature vector points into the chosen domain $\Omega$ wherever it is nonzero. Then there is a unique smooth complete embedded $E$-minimizing self-expanding hypersurface $M\subset\Omega$ with tangent cone $C$ at infinity. Its mean curvature vector is everywhere nonzero and points into $\Omega$, and
\[
 \{\sqrt t\,M:t>0\}
\]
is a foliation of $\Omega$ moving by mean curvature flow. The normalization is $\vec H=X^\perp/2$.

This is an existence theorem for a family determined by the cone and the chosen side, not an explicit parametrization. The uniqueness just stated includes the chosen-side and minimizing hypotheses. Its topology, symmetry, and number of ends depend on the chosen cone. Source: \cite{CatDing20}, Theorem 1.1, printed p.~361, supplied PDF p.~3, and Theorem 6.12, printed pp.~402--403, PDF pp.~44--45. The broader minimizing-current existence theorem in that paper can have singularities and is not being substituted for this smooth theorem. The figure is a schematic of the cone, selected domain, and expanding leaves.
\end{example}

\chapter{Self-expanding planar tree networks}\label{sec:catalogue-expanding-networks}
\begin{example}[Regular expanding trees with prescribed rays]\label{ex:catalogue-expanding-networks}
A regular planar network consists of smooth branches meeting only at triple junctions, with angles $120$ degrees at every junction. A self-expanding network $\Sigma\subset\mathbb R^2$ satisfies
\[
 \vec k=\frac{X^\perp}{2}
\]
on each open branch; its time evolution is $\Sigma_t=\sqrt t\,\Sigma$, $t>0$, with the same angle condition. It is a network flow and is not a smooth immersed curve at its junctions. For finitely many distinct prescribed rays from the origin, connected tree-like expanders provide smoothings in this network sense. Existence for three rays, including uniqueness of the tree, and existence for more rays are recalled in \cite{CatINS19}, Section 7, printed p.~70, supplied PDF p.~32, with attribution there to Schn\"urer--Schulze and Mazzeo--S\'aez respectively.

The branches can be described as geodesics for the conformal metric
\[
 g=e^{|X|^2/2}(dX_1^2+dX_2^2),
\]
with the junction balance condition. Noncompact ends approach their assigned rays with Gaussian decay. Two regular expanders with the same asymptotic rays coincide if they are in the same topological class: the comparison deformation must retain distinct junctions and the prescribed asymptotic ends, as in the source definition. This does not assert uniqueness across different network topologies. See \cite{CatINS19}, Lemmas 4.1 and 4.3 and Theorem 4.4, printed pp.~49--51, supplied PDF pp.~11--13.

The source uses $\vec k=\psi^\perp$ and $\sqrt{2t}\,\psi$; our profile is $\Sigma=\sqrt2\,\psi$. Nontrivial networks arise from a geodesic boundary-value construction. The exact balanced special case
\[
 \bigcup_{j=0}^2\{r(\cos(2\pi j/3),\sin(2\pi j/3)):r\geq0\}
\]
is the stationary $Y$ network. It illustrates the junction condition.

The accompanying picture instead computes a nonstationary four-ended example, using the geodesic construction of \cite{CatMS07}, Sections 2--3. In Euclidean arclength $s$, write the tangent as $(\cos\theta,\sin\theta)$. Then
\[
 x'=\cos\theta,\qquad y'=\sin\theta,\qquad
 \theta'=\frac{-x\sin\theta+y\cos\theta}{2}.
\]
Shoot the upper-right branch from $(x,y,\theta)=(a,0,\pi/3)$ toward the ray of angle $\pi/4$. Numerically, $a\approx0.385405$ in our normalization. Reflections in the coordinate axes and the straight segment $[-a,a]\times\{0\}$ give the displayed tree, with triple junctions at $(\pm a,0)$ and asymptotic rays at angles $\pi/4+j\pi/2$, $j=0,1,2,3$. A quarter-turn gives the second tree. The two pictures are congruent by rotation but have different connections when the four asymptotic ends are held fixed. Both have four ends and no cycles. These are numerical ODE approximations, not closed formulas or a classification of all network expanders. The time control shows $\sqrt t\,\Sigma$, $t>0$; the $t=0$ view shows the nonregular initial cone only.
\end{example}

\chapter{Mountain-pass self-expanders}\label{sec:catalogue-mountain-pass-expanders}
\begin{example}[Mountain-pass self-expanders]\label{ex:catalogue-mountain-pass-expanders}
Let $2\leq n\leq6$, let $C\subset\mathbb R^{n+1}$ be a $C^3$-regular cone, and suppose two distinct strictly stable $C^2$-asymptotically conical self-expanders $\Gamma_-\preceq\Gamma_+$ approach $C$. The order means inclusion of their consistently oriented enclosed regions. Bernstein--Wang construct a further smooth properly embedded self-expander $\Gamma_0$ between them, meeting the open intervening region. The associated flow is $M_t=\sqrt t\,\Gamma_0$, $t>0$, with $\mathbf H=X^\perp/2$.

If the cone is generic in the precise sense that every expander asymptotic to it has no nontrivial Jacobi field fixing infinity, $\Gamma_0$ can be chosen unstable for expander weighted area. The theorem does not prescribe its genus, connectedness, symmetry or uniqueness. Completeness follows from proper embeddedness without boundary. Cone asymptotics give Euclidean area growth and hence finite Gaussian area and entropy; no numerical value is supplied.

Corollary 1.3 gives an open set of cones admitting at least two connected and one disconnected smoothing, beyond the rotational setting. Sources: \cite{CatBWMP20}, Theorem 1.1 and Corollaries 1.2--1.3, pp.~1--2. In higher dimensions the general theorem allows a singular set of codimension seven; only its smooth $2\leq n\leq6$ cases are catalogued here. The figure depicts ordering, not a computed profile or prescribed topology.
\end{example}

\chapter{Expanders of unbounded genus at a fixed cone}\label{sec:catalogue-high-genus-expanders}
\begin{example}[Expanders of unbounded genus at a fixed cone]\label{ex:catalogue-high-genus-expanders}
Hoffman--Mart\'in--White construct connected embedded self-expanders in $\mathbb R^3$ of arbitrarily large genus at a single suitable cone. Precisely, let $C=C_k$ satisfy the symmetry and three-component link hypotheses of their Theorem 1.1, with its spherical link sufficiently close to the equatorial plane as required in Theorem 1.2. For each odd prime $p$ and integer $a\geq0$, there are at least four noncongruent $G_k$-invariant expanders asymptotic to $C$, of genus $p^a k-1$.

These are complete proper noncompact surfaces with three conical ends; the end count follows from the three components of the link. They solve $\mathbf H=X^\perp/2$ and yield $\sqrt t\,E$ for $t>0$. They are obtained by degree theory and limiting constructions, not an explicit formula. No common convexity, stability, or numerical entropy value is asserted. Smooth conical ends imply finite Gaussian area and entropy.

The incoming shrinker $M_k$ has genus $2k-2$. Joining $\sqrt{-t}\,M_k$ for $t<0$, $C$ at $t=0$, and $\sqrt t\,E$ for $t>0$ gives the source's self-similar Brakke flow, smooth except at the spacetime origin. This is not a genus increase in a smooth flow across $t=0$. Source: \cite{CatHMWExpand26}, Theorems 1.1--1.2 and equation (1.5), pp.~2--3. The figure is an original topology model, not a numerical solution of the expander equation. The three ends wrap around the $z$-axis, rather than forming three separate tubes. Its displayed cone and small symmetry counts are illustrative; they are not asserted to be the cone $C_k$ of a constructed shrinker or to satisfy the theorem's quantitative hypotheses.

The two slice pairings follow Definition 3.1 and Figure 1 (pp.~9--10): in the reflection plane $\{y=0\}$, the middle end on the positive-$x$ side is joined to the upper end for type~1, and to the lower end for type~2, with no closed slice curves. Each type has both a small and a big expander. These terms refer to the local integral $\int_E\phi\,d\mathcal H^2$, below or above the threshold $3/2$, where $\phi$ is a fixed nonnegative radial function supported in $B(0,1/2)$ and normalized by $\int_{\{z=0\}}\phi\,d\mathcal H^2=1$ (Definition 4.1 and Theorem 4.3). They do not refer to the magnification of the picture. Theorem 5.1 (p.~12) gives all four classes, and Theorems 12.6--12.7 (pp.~39--40) retain genus and class when the boundary radius tends to infinity.
\end{example}

\chapter{Knotted self-expanders}\label{sec:catalogue-knotted-expanders}
\begin{example}[Knotted self-expanders]\label{ex:catalogue-knotted-expanders}
Mramor--Shengwen Wang construct complete properly embedded self-expanders in $\mathbb R^3$ with knotted topology. Their Theorem 3.1 starts with an embedded orientable surface bounding a handlebody and, with sufficiently many ends, produces an expander whose bounded-region topology realizes that handlebody. These surfaces are stable for the weighted functional $\int e^{|X|^2/4}\,d\mu$; this is expander weighted stability, not entropy stability.

A concrete case has genus one and three ends: inside large balls it is isotopic to the boundary of a tubular neighborhood of a trefoil knot with three disks removed and annular ends attached. The half-cylinders in this description specify topology; the constructed geometric ends are asymptotically conical. More generally, every genus $g\geq1$ occurs with sufficiently many ends, without a universal end count specified here.

The normalization is $\mathbf H=X^\perp/2$, with flow $M_t=\sqrt t\,M$ for $t>0$. The examples are noncompact, complete and proper. No explicit formula, uniqueness, or exact Gaussian-area or entropy value is provided; conical area growth gives finiteness. Source: \cite{CatMWKnotted26}, Theorems 1.1 and 3.1 and Corollary 3.2, pp.~2 and 8--10, arXiv v1 (November 2, 2026). This is a preprint construction, not independently verified here. The interactive picture shows a trefoil tube with three annular ends, with auxiliary views of the closed tube and its three punctures. The accompanying plate explains these topological operations; Figure 2 of the paper motivates the bridge description. All displayed shapes are schematics, not computed expanders.
\end{example}

\chapter{Joyce--Lee--Tsui Lagrangian expander necks}\label{sec:catalogue-jlt-expanders}
These expanders smooth two transverse Lagrangian planes in $\mathbb C^n$.
\begin{example}[Expanding Lagrangian necks]\label{ex:catalogue-jlt-expanders}
Fix $n\geq2$, $a_j>0$, and $\psi_j\in\mathbb R$. Set
\[
 P_\alpha(y)=\frac{e^{\alpha y^2}\prod_{j=1}^n(1+a_jy^2)-1}{y^2},
 \quad P_\alpha(0)=\alpha+\sum_ja_j,\quad
 \phi_j(y)=\psi_j+\int_0^y\frac{dt}{(a_j^{-1}+t^2)\sqrt{P_\alpha(t)}}.
\]
With $\alpha=1/2$, the exact integral parametrization
\[
 X(x,y)=\left(x_j\sqrt{a_j^{-1}+y^2}\,e^{i\phi_j(y)}\right)_{j=1}^n,
 \qquad (x,y)\in\mathbb S^{n-1}\times\mathbb R,
\]
satisfies $\vec H=X^\perp/2$ and gives the flow $\sqrt t\,X$, $t>0$.
By \cite{CatJoyceLeeTsui2010}, Theorems~C--D, its image is a proper embedded Lagrangian neck. It is complete, has dimension and codimension $n$, and has two ends. For $n=2$ its topology is an annulus; generally it is $\mathbb S^{n-1}\times\mathbb R$.

Write $\bar\phi_j=\phi_j(+\infty)-\psi_j$. The limiting planes have coordinate phases $\psi_j\pm\bar\phi_j$, with $0<2\sum_j\bar\phi_j<\pi$. Every admissible angle tuple determines unique positive $a_j$ at fixed $\alpha$ within this family. The neck is zero-Maslov and almost calibrated. Equal $a_j$ and equal $\psi_j$ give rotational symmetry under $\mathrm{SO}(n)$.

For $n\geq3$, the stronger uniqueness theorem in \cite{CatImagiJoyceOliveira2016}, Theorem~1.2, requires an embedded, exact, closed-as-a-subset Lagrangian expander asymptotic to two transverse planes at rate $\rho<2$. These hypotheses must not be dropped. No general stability claim is attached. The figure projects the $n=2$, $a_1=a_2=1$, $\psi_1=\psi_2=0$ member, evaluating its integrals numerically.
\end{example}

\chapter{Two-convex Lagrangian graphical expanders}\label{sec:catalogue-ttw-expanders}
Two-convexity here is a condition on a potential's Hessian, not hypersurface mean convexity.
\begin{example}[Expanders from homogeneous potentials]\label{ex:catalogue-ttw-expanders}
Let $n\geq2$. Suppose $U_0:\mathbb R^n\to\mathbb R$ is homogeneous of degree two, is $C^2$ away from the origin, and its Hessian eigenvalues obey
\[
 \lambda_i+\lambda_j\geq0,\qquad 1+\lambda_i\lambda_j\geq0\quad(i\ne j),
 \qquad \prod_{i=1}^n(1+\lambda_i^2)^{-1/2}\geq\varepsilon_1>0.
\]
Proposition~6.5 of \cite{CatTsaiTsuiWang2025} produces the unique solution in the stated continuous-initial-data, smooth-positive-time class of
\[
 u_t=\sum_{i=1}^n\arctan\lambda_i(D^2u),\qquad
 u(x,0)=U_0(x),\qquad u(x,t)=t\,u(x/\sqrt t,1).
\]
Then $L_t=\{(x,Du(x,t)):x\in\mathbb R^n\}\subset\mathbb C^n$, $t>0$, is self-expanding; $L_1$ satisfies $\vec H=X^\perp/2$. These are complete proper embedded graphs of dimension and codimension $n$, homeomorphic to $\mathbb R^n$, with one end. Hessian bounds persist and higher derivatives obey the estimates in Theorem~5.6. The initial cone is prescribed through the potential; no stronger asymptotic rate is asserted here.

For a concrete nonquadratic seed, take $U_0(0)=0$ and
\[
 U_0(x)=-\frac{x_1^2}{20}+\frac14\sum_{j=2}^nx_j^2
          +10^{-4}\frac{x_1^4}{|x|^2}\quad(x\ne0).
\]
A direct differentiation bounds the operator norm of $D^2(x_1^4/|x|^2)$ by $38$. Thus one Hessian eigenvalue lies within $0.0038$ of $-0.1$, and the others within $0.0038$ of $0.5$. Both pair inequalities are strictly positive. This verifies the theorem's hypotheses while keeping the potential nonconvex. This formula describes initial data, not the positive-time expander. The illustration shows its exact initial gradient cone and marks the subsequent PDE smoothing as a construction step.
\end{example}

\chapter{Lee--Wang embedded quadric expanders}\label{sec:catalogue-lw-quadric-expanders}
\begin{example}[The expanding mixed-sign quadric family]\label{ex:catalogue-lw-quadric-expanders}
For $n\geq3$, $n/2<k<n$, and $S=2k-n$, set
\[
 Q_-=\left\{x\in\mathbb R^n:\sum_{j\leq k}x_j^2-\sum_{j>k}x_j^2=-2S\right\},\qquad
 X(x,s)=(e^{is}x_1,\ldots,e^{is}x_k,e^{-is}x_{k+1},\ldots,e^{-is}x_n).
\]
On $(Q_-\times\mathbb R)/((x,s)\sim(-x,s+\pi))$ this is a complete proper embedded Lagrangian in $\mathbb C^n$, of dimension and codimension $n$. It satisfies $\vec H=X^\perp/2$, and $L_t=\sqrt t\,X(Q_-\times\mathbb R)$ for $t>0$. This is an exact parametrization subject to one quadric constraint, not an ODE approximation. See \cite{CatLeeWang2010}, Propositions~2.1--2.2, with weights $+1$ and $-1$.

The fiber is $\mathbb R^k\times S^{n-k-1}$. The mapping-torus quotient is connected and has one end. If $n-k=1$, its two fiber components are exchanged by the gluing, and the total manifold is diffeomorphic to $\mathbb R^k\times S^1$, hence orientable. For $n-k\geq2$ the fiber is connected, the antipodal monodromy must be retained, and the total manifold is orientable exactly when $n$ is even. In particular, the $n=3,k=2$ case is orientable: one must pass around the base twice to return to the same component. These topology statements are obtained from the quotient directly; the product and parity shorthand in the source's Proposition~2.2 requires this disconnected-fiber qualification.

The family is noncompact and Hamiltonian stationary, with $SO(k)\times SO(n-k)$ and circle symmetries, and is asymptotic to the same cone as Example~\ref{ex:catalogue-lw-quadric-shrinkers}. Its planar-coordinate quadric description gives polynomial volume growth and finite centered Gaussian area; this entry records no entropy value or stability result. No surface genus is assigned to these $n\geq3$ examples. The source's weak continuation at $t=0$ is a separate varifold statement; the smooth time interval here is $t>0$.

For the illustrated case $n=3$, $k=2$, the full expander has the global parametrization
\[
 F_t(a,b,s)=\left(e^{is}a,e^{is}b,e^{-is}\sqrt{a^2+b^2+2t}\right),
 \qquad (a,b)\in\mathbb R^2,\quad s\in\mathbb R/(2\pi\mathbb Z),\quad t>0.
\]
This parametrization is obtained directly from the unit-weight construction in \cite{CatLeeWang2010}, Theorem~1.1 and Section~2.2: one takes the positive component of the real quadric and a full circle of phases. The real slice, obtained at $s=0$ and $s=\pi$, consists of the two graphs
\[
 x_3=\pm\sqrt{r^2+2t},\qquad r=\sqrt{x_1^2+x_2^2}\geq0.
\]
Their tips are at $x_3=\pm\sqrt{2t}$. The figure compares $t=1/4,1,4$ with the singular cone limit $t=0$, using the same spatial scale and the fixed display cutoff $r\leq3.2$. The complete slices are unbounded. These are exact two-dimensional slices of a three-dimensional flow in $\mathbb R^6$; they are not themselves mean curvature flows in $\mathbb R^3$.

The circle direction explains why two real sheets do not mean two components of the expander. For a point $x=(x_1,x_2,x_3)$ on the positive real sheet, the curve
\[
 \gamma_x(s)=\left(e^{is}x_1,e^{is}x_2,e^{-is}x_3\right)
 =\cos(s)x+\sin(s)i(x_1,x_2,-x_3),\qquad 0\leq s\leq\pi,
\]
lies in $L_t$ and joins $x$ to $-x$ on the other real sheet. It is a semicircle in a real two-plane in $\mathbb R^6$, since the two spanning vectors are orthogonal and have equal length. The circle-phase control follows this path at fixed time; only the time control changes the flow.
\end{example}
